% Einheit 5 von 8 – Mammutbäume (bank/zufallsexperimente-und-pfadregeln/e5.jsonl, zusammenbau v0.1)
%% TODO Zeitmarke Einheit 5: Marken-Zeile nicht lesbar
%% TODO Zweigzeile Teil 1: was hier gelernt wird (Satz in Schülersprache aus dem Einheitstitel) – Einheit 5
\einheitenkopf[e5]{Einheit 5 von 8 · Mammutbäume}
\zweigzeile{Abitur GK}
\setcounter{aufgabe}{22}

%% TODO Titel: Ich-kann-Satz fehlt in der Bank (Platzhalter: Kettenname) – Nr. 23
%% TODO Anweisung: ein Satz über der ersten Teilaufgabe fehlt in der Bank – Nr. 23
\begin{aufgabe}{Mammutbäume}
%% TODO \rechenplatz{n}: Schrittzahl des Grundfalls fehlt in der Bank (nur bei fester Schreibform, 2.3 b)
\begin{teile}
\teil Ein Würfel wird dreimal geworfen: „genau eine Sechs“. Ein Pfad oder mehrere? Wie viele?\\ \kreuz{genau ein Pfad}\\ \kreuz{mehrere Pfade mit denselben Faktoren}
\teil Ein Glücksrad liefert einen Gewinn mit $0{,}2$; es wird dreimal gedreht. P(genau ein Gewinn)? P = \leerfeld
\teil Ein Glücksrad zeigt Rot mit $0{,}5$, Gelb mit $0{,}3$ und Grün mit $0{,}2$; es wird dreimal gedreht. P(drei verschiedene Farben)? P = \leerfeld
\teil Ein Glücksrad zeigt 1 mit $0{,}5$, 2 mit $0{,}3$ und 3 mit $0{,}2$; es wird dreimal gedreht. P(Summe 6)? P = \leerfeld
\teil Ein Versuch gelingt mit $0{,}6$. Sechs Versuche. P(genau drei Erfolge, und zwar unmittelbar hintereinander)? P = \leerfeld
\teil Eine Urne enthält 4 rote und 3 weiße Kugeln; drei werden ohne Zurücklegen gezogen. P(mehr rote als weiße)? P = \leerfeld
\teil Von 10 Losen sind 4 Gewinne; man zieht 3. P(genau 2 Gewinne) mit Binomialkoeffizienten? P = \leerfeld
\teil Ein Glücksrad zeigt 1 mit $\frac{1}{2}$, 2 mit $\frac{1}{3}$ und 3 mit $\frac{1}{6}$; es wird n-mal gedreht ($n > 2$). Term für P(das Produkt aller Zahlen ist 3)? Werte für $n = 4$ aus. \hfill (Abitur 2025 LK) \\ P = \leerfeld
\end{teile}
\end{aufgabe}

%% TODO Titel: Ich-kann-Satz fehlt in der Bank (Platzhalter: Kettenname) – Nr. 24
%% TODO Anweisung: ein Satz über der ersten Teilaufgabe fehlt in der Bank – Nr. 24
\begin{aufgabe}{Mammutbäume – Fehler finden, Begründen, Anwendung, Darstellungswechsel}
\begin{teile}
\teil Ein Würfel wird dreimal geworfen. Max: P(genau zwei Sechsen) $= \left(\frac{1}{6}\right)^2 \cdot \frac{5}{6} = \frac{5}{216}$. Finde den Fehler.
\teil Begründe: Pfade mit denselben Sorten in anderer Reihenfolge haben dieselbe Wahrscheinlichkeit – auch ohne Zurücklegen.
\teil Eine Bäckerei verkauft Körnerbrötchen mit $0{,}4$, Weizen mit $0{,}35$ und Roggen mit $0{,}25$. Drei Kunden kaufen je ein Brötchen. P(jede Sorte genau einmal)? P = \leerfeld
\teil Eine Urne enthält 4 rote und 5 blaue Kugeln; drei ohne Zurücklegen, genau eine rote. Schreibe statt des Pfadterms den Lotto-Bruch und berechne.
\end{teile}
\end{aufgabe}
